
\documentclass[a4paper,10pt]{article}
\usepackage[margin=1.85cm]{geometry}
\usepackage[T1]{fontenc}
\usepackage[utf8]{inputenc}
\usepackage[italian]{babel}
\usepackage{lmodern,booktabs,tabularx,xcolor,listings,tikz,hyperref}
\usetikzlibrary{arrows.meta,positioning}
\definecolor{ink}{HTML}{163047}
\definecolor{pale}{HTML}{EEF4F7}
\definecolor{warn}{HTML}{934618}
\lstset{basicstyle=\ttfamily\fontsize{7.8}{9.1}\selectfont,breaklines=true,
breakatwhitespace=false,columns=fullflexible,keepspaces=true,showstringspaces=false,
backgroundcolor=\color{pale},frame=single,rulecolor=\color{pale},
xleftmargin=4pt,xrightmargin=4pt,aboveskip=6pt,belowskip=6pt}
\setlength{\parindent}{0pt}\setlength{\parskip}{5pt}
\newcommand{\tagline}[1]{{\small\color{ink}\textbf{#1}}\par}
\begin{document}
\tagline{INTEGRAZIONE AL CONFRONTO SUL TEST | 27 settembre 2026}
{\LARGE\bfseries Qualità e correttezza\\dei fatti verificabili}\par
Analisi dei registri delle 90 decisioni LLM + RAG e delle 90 decisioni LLM senza RAG già raccolte.
Questo documento integra il confronto esistente senza modificarlo e senza eseguire nuovi Test.
\section*{Che cosa è verificato}
La risposta JSON separa \texttt{facts} da \texttt{hypothesis}. I fatti sono asserzioni scelte da quattro
tipi predefiniti: il programma confronta campi precisi con i dati disponibili. Non valuta la verità di un testo libero.

\textbf{E1, \ldots, E5 identificano gli esempi recuperati, non i fatti.}
Ogni fatto con riferimento indica un alias locale. La posizione dell'esempio nel prompt lo collega
a un record train preciso. E1 può quindi indicare circuiti diversi in richieste diverse.
\begin{center}
\begin{tikzpicture}[node distance=5mm,every node/.style={draw=ink,rounded corners,
align=center,font=\small,fill=pale,minimum height=13mm},>=Latex]
\node(a){Fatto nel JSON\\\texttt{example\_id: Ei}};
\node(b)[right=of a]{Alias Ei\\nel prompt};
\node(c)[right=of b]{Record train\\\texttt{rag\_id}};
\node(d)[right=of c]{Campo del record\\confrontato};
\draw[->,thick](a)--(b);\draw[->,thick](b)--(c);\draw[->,thick](c)--(d);
\end{tikzpicture}\end{center}
\section*{Risultati su tutti i casi}

\begin{tabularx}{\linewidth}{Xrr}\toprule Misura & LLM + RAG & Senza RAG\\\midrule
Risposte finali con tutti i fatti validi & 73/90 & 90/90\\
Risposte finali con almeno un fatto non valido & 17/90 & 0/90\\
Fatti validi nelle risposte finali & 163/180 & 90/90\\
Risposte valide già al primo tentativo & 71/90 & 90/90\\
Risposte complete, incluse le correzioni & 126 & 90\\
\bottomrule\end{tabularx}\par
Nel RAG, 95 fatti finali riferiti agli esempi risultano tutti coerenti con i campi citati. I 17 fatti non validi sono sempre fatti sulla capacità del dispositivo con un riferimento Ei non consentito. La capacità risulta comunque sufficiente; l'errore è nell'attribuzione dell'evidenza. Dopo tre tentativi queste scelte sono state accettate dalla regola prevista, conservando lo stato di mancata verifica.\par
\textbf{Il 90/90 senza RAG riguarda solo la capacità in qubit.} È un controllo più semplice: non dimostra spiegazioni migliori né prestazioni di compilazione superiori.\par
\textbf{Le ipotesi non sono verificate sul piano dei contenuti e non vanno prese come affermazioni vere.} I controlli di formato, lunghezza o scelta ammessa non certificano ciò che il testo sostiene.\par
\newpage
\tagline{CAMPIONE 1 / 3 | selezione illustrativa}
\section*{Due fatti sostenuti dallo stesso esempio}
\textbf{Circuito Test:} \texttt{groundstate\_small\_indep\_tket\_4}. Risposta finale, tentativo 1.\par
\textbf{JSON originale dell'LLM.} Valori integrali; sola impaginazione modificata.\par
\begin{lstlisting}
{
  "selected_device": "ibm_falcon_127",
  "config_id": "o2_dense_sabre",
  "facts": [
    {
      "assertion": "selected_device_matches_example",
      "example_id": "E1"
    },
    {
      "assertion": "selected_pair_among_reported_best",
      "example_id": "E1"
    }
  ],
  "hypothesis": "This pair is selected because the device and configuration appear together as top performers in the historical record for a similar four-qubit circuit. The hypothesis suggests that applying this specific setup to the current circuit may yield comparable results, though no performance guarantee exists for this new instance."
}
\end{lstlisting}

\textbf{Collegamento all'evidenza train}
\begin{center}\begin{tikzpicture}[>=Latex,every node/.style={draw=ink,fill=pale,rounded corners,font=\small,align=center}]\node(a){\texttt{example\_id: E1}};\node(b)[right=8mm of a]{E1 nel prompt};\node(c)[right=8mm of b]{\texttt{qaoa\_indep\_qiskit\_4}\\esempio train};\draw[->,thick](a)--(b);\draw[->,thick](b)--(c);\end{tikzpicture}\end{center}
{\footnotesize Record completo: \texttt{rag\_41b126cd2f33c412dbe984b3b8ce1bb72143184b5521243a890e29f7b64b74b5}}\par
\textbf{Estratto del record (prima riga della graduatoria)}
\begin{lstlisting}
{
  "id": "E1",
  "num_qubits": 4,
  "selected_device": "ibm_falcon_127",
  "top_configurations": [
    {
      "rank": 1,
      "device_id": "ibm_falcon_127",
      "config_id": "o2_dense_sabre",
      "median_score": 0.9822170185,
      "tied_score_config_ids": [
        "o2_dense_sabre",
        "o3_dense_sabre"
      ]
    }
  ]
}
\end{lstlisting}

\textbf{Esito del controllo.} Il dispositivo scelto coincide con quello selezionato nel record. Anche la coppia Falcon 127 / o2\_dense\_sabre compare nella graduatoria mostrata, al primo posto. È presente una parità con o3\_dense\_sabre: il fatto non certifica un vincitore unico. Il valore 0,9822170185 è la mediana storica di QAOA, non una misura del circuito groundstate corrente.\par
{\color{warn}\textbf{Limite dell'ipotesi.}} L'ipotesi propone che il risultato si trasferisca al circuito nuovo. Questo passaggio non è verificato.\par
{\small Stati registrati, in ordine: \texttt{verified, verified}. Stato della risposta: \texttt{success}.}
\newpage
\tagline{CAMPIONE 2 / 3 | selezione illustrativa}
\section*{Coppia storica e uguaglianza dei qubit}
\textbf{Circuito Test:} \texttt{qpeexact\_indep\_qiskit\_2}. Risposta finale, tentativo 2.\par
\textbf{JSON originale dell'LLM.} Valori integrali; sola impaginazione modificata.\par
\begin{lstlisting}
{
  "selected_device": "ibm_falcon_127",
  "config_id": "o2_default_default",
  "facts": [
    {
      "assertion": "selected_pair_among_reported_best",
      "example_id": "E1"
    },
    {
      "assertion": "same_qubit_count_as_example",
      "example_id": "E1"
    }
  ],
  "hypothesis": "This pair is selected because it matches the top-performing configuration for a similar 2-qubit circuit on the same hardware, and the qubit count is identical to the reference example."
}
\end{lstlisting}

\textbf{Collegamento all'evidenza train}
\begin{center}\begin{tikzpicture}[>=Latex,every node/.style={draw=ink,fill=pale,rounded corners,font=\small,align=center}]\node(a){\texttt{example\_id: E1}};\node(b)[right=8mm of a]{E1 nel prompt};\node(c)[right=8mm of b]{\texttt{dj\_indep\_qiskit\_2}\\esempio train};\draw[->,thick](a)--(b);\draw[->,thick](b)--(c);\end{tikzpicture}\end{center}
{\footnotesize Record completo: \texttt{rag\_251477152075954d06ead6d2f929526b197b4e3f572f8c9ae5f01bf4481156cd}}\par
\textbf{Estratto del record (prima riga della graduatoria)}
\begin{lstlisting}
{
  "id": "E1",
  "num_qubits": 2,
  "selected_device": "ibm_falcon_127",
  "top_configurations": [
    {
      "rank": 1,
      "device_id": "ibm_falcon_127",
      "config_id": "o2_default_default",
      "median_score": 0.9994380656,
      "tied_score_config_ids": [
        "o2_default_default",
        "o3_default_default"
      ]
    }
  ]
}
\end{lstlisting}

\textbf{Esito del controllo.} La coppia scelta è presente nelle righe mostrate del record Deutsch--Jozsa. Il circuito corrente e l'esempio hanno entrambi 2 qubit: il secondo fatto è quindi vero. L'uguaglianza del numero di qubit non prova che topologia, sequenza delle porte o risposta alla compilazione siano simili.\par
{\color{warn}\textbf{Limite dell'ipotesi.}} La formulazione \emph{similar 2-qubit circuit} è una valutazione libera. Il validatore controlla l'uguaglianza 2 = 2, non la similarità strutturale né il trasferimento delle prestazioni.\par
{\small Stati registrati, in ordine: \texttt{verified, verified}. Stato della risposta: \texttt{success}.}
\newpage
\tagline{CAMPIONE 3 / 3 | fatto non valido}
\section*{Fallimento del controllo: evidenza non ammessa}
\textbf{Circuito Test:} \texttt{qpeexact\_indep\_qiskit\_13}.
Primo tentativo, seguito da una richiesta di correzione: non è la risposta finale.\par
\textbf{JSON originale dell'LLM.} Contenuto integrale; sola impaginazione modificata.\par
\begin{lstlisting}
{
  "selected_device": "quantinuum_h2_56",
  "config_id": "o2_default_default",
  "facts": [
    {
      "assertion": "selected_device_matches_example",
      "example_id": "E1"
    },
    {
      "assertion": "selected_device_has_enough_qubits",
      "example_id": "E1"
    }
  ],
  "hypothesis": "This pair is selected because the device matches the top choice in the historical examples and possesses sufficient qubits for the current circuit. The hypothesis is a proposal, not a verified fact: do not claim that performance on the current circuit is known or guaranteed. Historical results belong only to their original circuits."
}
\end{lstlisting}

\textbf{E1 esiste, ma non sostiene entrambi i fatti.} Nel prompt corrisponde a questo esempio train:
\begin{lstlisting}
{
  "id": "E1",
  "circuit_id": "graphstate_indep_qiskit_14",
  "num_qubits": 14,
  "selected_device": "quantinuum_h2_56"
}
\end{lstlisting}

\begin{center}
\begin{tikzpicture}[>=Latex,every node/.style={rounded corners,draw=ink,fill=pale,
font=\small,align=center,minimum height=10mm}]
\node(a){Fatto 1: stesso dispositivo};
\node(b)[right=7mm of a]{E1: Quantinuum H2-56};
\node(c)[right=7mm of b]{\texttt{verified}};
\node(d)[below=5mm of a]{Fatto 2: capacità sufficiente};
\node(e)[below=5mm of b,draw=warn]{E1: fonte non ammessa};
\node(f)[below=5mm of c,draw=warn]{\texttt{unsupported}};
\draw[->,thick](a)--(b);\draw[->,thick](b)--(c);
\draw[->,thick,draw=warn](d)--(e);\draw[->,thick,draw=warn](e)--(f);
\end{tikzpicture}
\end{center}
\textbf{Errore registrato dal validatore (estratto).} L'indice 1 identifica il secondo fatto.
\begin{lstlisting}
{
  "facts_status": "unverified",
  "fact_checks": [{"index": 1, "result": "unsupported"}],
  "issues": [{"code": "FACT_NOT_SUPPORTED", "path": "$.facts[1]"}]
}
\end{lstlisting}
\textbf{Perché fallisce.} L'asserzione \texttt{selected\_device\_has\_enough\_qubits}
deve usare la capacità del catalogo hardware e non deve contenere \texttt{example\_id}.
Il modello cita invece E1. I 14 qubit dell'esempio descrivono quel circuito train,
non la capacità del dispositivo. Il dato pertinente è 56 qubit disponibili per i 13 richiesti.\par
\textbf{Non è una capacità falsa, ma un fatto non valido nel contratto di risposta.}
La presenza di un esempio reale non basta: deve essere la fonte ammessa per l'asserzione.
Il JSON sopra resta quello prodotto dal modello, senza correzioni.\par
\textbf{Che cosa succede dopo.} Il controllo richiede una correzione. L'errore persiste fino al
tentativo 3; la coppia viene infine accettata come \texttt{accepted\_with\_unverified\_facts}.
La compilazione riesce: il fallimento illustrato riguarda la verifica del fatto.\par
{\color{warn}\textbf{Ipotesi non verificata.}} La prosa non viene validata semanticamente,
anche quando ripete le istruzioni del prompt, e non va letta come un'affermazione vera.\par
\newpage\tagline{PORTATA DEI RISULTATI E RIPRODUCIBILITÀ}
\section*{Correttezza dei fatti e qualità della spiegazione}
\begin{tabularx}{\linewidth}{Xrr}\toprule Tipo di fatto finale RAG & Validi & Totali\\\midrule
Stesso dispositivo dell'esempio & 88 & 88\\
Coppia tra le configurazioni mostrate & 5 & 5\\
Stesso numero di qubit & 2 & 2\\
Capacità in qubit sufficiente & 68 & 85\\
\bottomrule\end{tabularx}\par
Il sistema rende controllabile una parte limitata della motivazione. La maggior parte dei fatti storici riguarda solo il dispositivo: 88 asserzioni. Soltanto 5 riguardano esplicitamente la coppia dispositivo/configurazione. Le 2 uguaglianze dei qubit descrivono la dimensione, senza dimostrare una somiglianza di comportamento.\par
\textbf{Fatto corretto non significa scelta ottima.} La presenza di una coppia nelle prime righe vale per quel circuito, quei Target sintetici e quella procedura. La regola sulla coppia accetta le righe mostrate e le parità dichiarate; non implica sempre primo posto né superiorità sul circuito corrente. La capacità in qubit è solo una condizione necessaria.\par
\textbf{Ipotesi: nessuna verifica semantica.} Il sistema non confronta la prosa con esperimenti, conoscenze esterne o una prova logica. Non certifica promesse di qualità, spiegazioni causali, similarità dei circuiti o bontà della configurazione. Anche una risposta con tutti i fatti validi mantiene \texttt{explanation\_fully\_verified=false}. Le tre osservazioni sulla prosa nelle pagine precedenti sono un commento qualitativo manuale, non una misura automatica né una valutazione sistematica di tutte le 90 ipotesi.\par
\section*{Come è stata svolta l'analisi}
Abbiamo riletto tutti i JSON delle risposte e dei controlli, compresi i tentativi scartati. Il JSON del tentativo finale coincide con la decisione conservata. Le mappe Ei, le impronte della vista e il contenuto degli esempi sono stati confrontati con prompt, registri e Dataset train verificato. Le quattro regole sono state riapplicate, anche con confronti espliciti separati dagli esiti salvati. I risultati coincidono con i controlli registrati.\par
Le 126 risposte RAG contengono complessivamente 252 fatti: 199 validi e 53 non validi. Questo denominatore include le correzioni e non va confuso con i 180 fatti finali. I campioni sono scelti per illustrare tre casi diversi; i conteggi complessivi provengono dall'intero insieme dei 90 circuiti.\par
\section*{Fonti locali e ricostruzione}
{\small\raggedright Generatore: \texttt{archivio/valutazione/test/report/analizza\_fatti.py}. Dati: \texttt{test/risultati/llm\_rag} e \texttt{llm\_senza\_rag}. Per ogni circuito: \texttt{prompt.json}, \texttt{encoding.json}, \texttt{decision\_validation.json}, \texttt{decision.json} e \texttt{attempt\_*/call/response\_raw.json}.\par}
Il file \texttt{audit\_fatti.json} conserva conteggi, controlli su tutti i tentativi, i tre JSON completi (due finali e un primo tentativo non valido), esempi train integrali e impronte SHA-256 delle fonti. La rianalisi non ripete le compilazioni storiche: la correttezza qui riguarda la coerenza con i dati forniti, non una verifica fisica della fedeltà.\par
\end{document}